% Generated draft: requires mathematical and source review.
\documentclass[11pt]{article}
\usepackage[margin=25mm]{geometry}
\usepackage{fontspec}
\setmainfont{Latin Modern Roman}
\usepackage{amsmath,amssymb,mathtools,unicode-math}
\setmathfont{Latin Modern Math}
\usepackage{enumitem,needspace,xurl,hyperref}
\setlength{\parindent}{0pt}
\setlength{\parskip}{5pt}
\setlength{\emergencystretch}{3em}
\newcommand{\N}{\mathbb{N}}
\newcommand{\Z}{\mathbb{Z}}
\newcommand{\Q}{\mathbb{Q}}
\newcommand{\R}{\mathbb{R}}
\newcommand{\C}{\mathbb{C}}
\newcommand{\F}{\mathbb{F}}
\newcommand{\A}{\mathbb{A}}
\newcommand{\PP}{\mathbb{P}}
\newcommand{\E}{\mathbb{E}}
\newcommand{\eps}{\varepsilon}
\newcommand{\dd}{\,\mathrm{d}}
\begin{document}
\section*{UnsolvedMath ID: 20002888}
\textbf{Source problem code:} catalogue entry
\subsection{Definitions and mathematical statement}
Bernhard Lamel asks about unbounded finite jet-determination order for automorphism germs. Let \(M\subset\C^n\) be a connected generic real-analytic submanifold, holomorphically nondegenerate and everywhere of finite type. Generic means \(T_pM+iT_pM=T_p\C^n\); finite type is understood in the CR bracket-generating sense. Holomorphic nondegeneracy excludes a nonzero holomorphic vector field complex-tangent to \(M\). Let \(\operatorname{Aut}(M,p)\) consist of germs of biholomorphic maps fixing \(p\) and preserving \(M\). Let \(G_p^k\) be the group of invertible \(k\)-jets of such ambient germs, and set
\[
 k(p)=\min\{k\ge1:f\mapsto j_p^kf
          \text{ is injective on }\operatorname{Aut}(M,p)\}.
\]
The source states finite determination at each point under its hypotheses. It asks whether an \(M\subset\C^2\) can have points \(p_j\) with \(k(p_j)\to\infty\), and comments that higher-dimensional examples exist. It then weakens finite type to generically finite type and asks whether such points can converge to a point of \(M\). In that variant, finite values \(k(p_j)\) must be required if the intended question is unbounded finite order, rather than simply assigning \(k=\infty\) at a point with no finite determination.

In \(\C^2\), the intended positive-CR-dimensional finite-type generic setting is a real hypersurface. The aid attributes to Ebenfelt--Lamel--Zaitsev the bound
\[
 k(p)\le2
\]
for real-analytic finite-type hypersurface germs in \(\C^2\), giving a negative answer to the everywhere-finite-type question. This scope should not be silently replaced by another definition of finite type for totally real submanifolds.

The bound is reported sharp. On the Heisenberg hypersurface \(\operatorname{Im}w=|z|^2\), the germs
\[
 F_t(z,w)=\left(\frac z{1+tw},\frac w{1+tw}\right),\qquad t\in\R,
\]
preserve the hypersurface and have identity first jet at the origin; \(t\ne0\) gives a nonidentity germ. Thus first jets do not suffice there.

For the generically finite-type accumulation question, the aid gives necessary restrictions rather than an example or nonexistence theorem. If \(p_j\to p\in M\) and finite \(k(p_j)\to\infty\), then eventually \(p_j\) lie in the infinite-type locus, since finite-type points have order at most two. The limit \(p\) is also infinite type because finite type is open. Jet order is invariant under local biholomorphic equivalence of germs, so infinitely many inequivalent germs must occur. For arbitrarily large \(N\), some such germ has nontrivial isotropy with
\[
 f\ne\mathrm{id},\qquad j_{p_j}^Nf=j_{p_j}^N\mathrm{id}.
\]
These are increasingly deep identity-jet kernels, not merely large values of a finite-type invariant.

The aid also records an analytic obstruction to one standard higher-dimensional construction. A fixed holomorphic function \(h\) on a connected neighborhood cannot have orders of vanishing tending to infinity at points converging inside that neighborhood unless \(h\equiv0\): every fixed derivative then vanishes at the limit. Thus a nonzero single-function shear/interpolation construction using zeros of increasing multiplicity at escaping nodes cannot simply replace those nodes by accumulating ones. This rules out that construction strategy, not every possible CR accumulation example. The original generically finite-type existence question remains unresolved in the supplied attempt.
\subsection{Short English statement}
The aid reports a sharp two-jet bound for finite-type real-analytic hypersurfaces in \(\C^2\), excluding unbounded orders in the first question. An accumulating example in the weaker generically finite-type setting would have to use infinite-type, increasingly inequivalent germs with deep identity-jet kernels. A fixed holomorphic interpolation function cannot supply increasing vanishing orders at accumulating interior points.
\subsection{Sources}
\begin{itemize}
\item Local UnsolvedMath registry, permanent ID 20002888, catalogue entry; canonical \texttt{aim-several-complex-variables-notes.json}, notes[45].
\item AIM \emph{Complexity of mappings in CR geometry}, Problem 6, proposed by Bernhard Lamel, \url{https://aimath.org/WWN/complexitycrmap/complexitycrmap.pdf}. The supplied verification reports reconstruction from pages 2--3, removing a running header inside ``finite type'' and restoring the jet notation; this draft does not independently inspect the PDF.
\item The aid attributes the sharp two-jet theorem to Ebenfelt, Lamel and Zaitsev, without a title, year, theorem number or bibliographic URL. The source's higher-dimensional-example comment likewise has no reference here.
\item Machine-generated research attempt 1, partial-result/partially-solved labels, reported rigor proved and low novelty confidence. Its accumulation reduction and construction obstruction are not peer reviewed and require independent verification.
\end{itemize}
\end{document}
